\section{引言}

\subsection{研究背景}
果树种植是全球农业生产的重要组成部分。准确的产量估算使农民能够优化收获调度、高效配置劳动力资源、规划仓储物流运作，并做出明智的市场定价决策~\citep{Maheswari_2021}。在精准农业中，及时准确的产量信息对于数据驱动的果园管理至关重要。

传统产量估算方法主要依赖人工视觉评估和统计采样。这些方法劳动强度大、耗时长，且存在显著的估算误差。研究表明，人工估算误差通常在30\%至50\%之间，主要原因是观察者经验差异、采样偏差以及对被遮挡果实的系统性低估~\citep{Sa_2016}。主观判断的依赖性造成了固有的不确定性，这些不确定性会传导到收获计划和市场预测中。

\subsection{消费级LiDAR技术}
消费电子产品中嵌入式固态LiDAR传感器的出现是一项重要的技术进步。苹果于2020年在iPad Pro中引入LiDAR传感器，随后扩展到iPhone 12 Pro及后续机型。该传感器采用间接飞行时间（iToF）技术，通过测量发射的调制红外激光脉冲与其反射信号之间的相位差来计算深度信息。

消费级LiDAR传感器的技术特性包括：有效测距范围约0.1至5.0米、每帧约30,000个点云的60 FPS输出、3米距离处约$\pm$5mm的深度精度。LiDAR与RGB相机和惯性测量单元（IMU）的集成使得带颜色的点云生成成为可能，设备价格区间为700至1500美元，远低于成本数万美元的专业地面激光扫描（TLS）系统。

\subsection{研究动机与目标}
尽管消费级LiDAR已在建筑扫描、室内设计和增强现实等领域得到应用，但其在农业产量估算中的利用仍是一个新兴研究方向。现有的基于LiDAR的果实检测研究主要使用昂贵的专业设备，而关于消费级LiDAR在果园环境中可行性的系统评估在文献中仍较为缺乏。

本研究旨在通过设计和实现使用消费级LiDAR设备的完整果实检测和产量估算系统来解决这一差距。具体目标包括：（1）评估消费级LiDAR在果园环境中的技术可行性；（2）设计和实现结合2D视觉检测与3D点云处理的多模态融合系统；（3）识别系统局限性并提出改进方向；（4）为后续真实果园数据的田间实验提供框架和基准。

\subsection{论文结构}
本文其余部分安排如下。第\ref{sec:related_work}节回顾果实检测方法、点云聚类算法和多模态融合方法的相关工作。第\ref{sec:methodology}节介绍系统架构和详细算法设计。第\ref{sec:experiments}节描述实验设置和初步结果。第\ref{sec:conclusion}节讨论系统局限性和未来工作。